\usepackage[a5paper,top=20mm,bottom=22mm,inner=21mm,outer=17mm,headheight=14pt]{geometry}
\usepackage{fontspec}
\usepackage{amsmath,mathtools}
\usepackage{unicode-math}
\setmainfont{texgyrepagella-regular.otf}[
  BoldFont=texgyrepagella-bold.otf,
  ItalicFont=texgyrepagella-italic.otf,
  BoldItalicFont=texgyrepagella-bolditalic.otf
]
\setsansfont{texgyreheros-regular.otf}[
  BoldFont=texgyreheros-bold.otf,
  ItalicFont=texgyreheros-italic.otf
]
\setmonofont{lmmono10-regular.otf}
\setmathfont{texgyrepagella-math.otf}

\usepackage{microtype}
\usepackage{setspace}
\setstretch{1.28}
\setlength{\parindent}{2em}
\setlength{\parskip}{0.15em}
\setlength{\emergencystretch}{3em}

\usepackage{booktabs,longtable,tabularx,array,multirow}
\usepackage{enumitem}
\setlist{nosep,leftmargin=2em}
\usepackage{xcolor}
\definecolor{deepblue}{HTML}{173B57}
\definecolor{warmgray}{HTML}{F3F1EC}
\definecolor{accent}{HTML}{9A5A36}
\usepackage[most]{tcolorbox}
\tcbset{enhanced,breakable,boxrule=.45pt,arc=1.2mm,left=1.5mm,right=1.5mm,top=1.2mm,bottom=1.2mm}
\newtcolorbox{learninggoals}{title=学习目标,colback=warmgray,colframe=deepblue,fonttitle=\bfseries}
\newtcolorbox{keypoint}{title=关键辨析,colback=blue!3,colframe=deepblue,fonttitle=\bfseries}
\newtcolorbox{workedexample}[1][]{title={完整算例：#1},colback=blue!3,colframe=deepblue,fonttitle=\bfseries}
\newtcolorbox{failurecheck}[1][]{title={失败诊断：#1},colback=red!2,colframe=accent,fonttitle=\bfseries}
\newtcolorbox{chapterreview}{title=本章小结,colback=warmgray,colframe=accent,fonttitle=\bfseries}
\newtcolorbox{selfcheck}{title=自测与研讨,colback=white,colframe=accent,fonttitle=\bfseries}

\usepackage{tikz}
\usetikzlibrary{arrows.meta,positioning,fit,calc,shapes.geometric}
\tikzset{
  causalnode/.style={draw=deepblue,thick,circle,minimum size=8mm,inner sep=1pt},
  causalbox/.style={draw=deepblue,rounded corners,align=center,minimum height=8mm,fill=blue!3},
  causalarrow/.style={-{Latex[length=2mm]},thick,deepblue}
}

\usepackage{csquotes}
\usepackage[backend=biber,style=gb7714-2015ay,gbpub=false,maxcitenames=2,mincitenames=1,sorting=nyt,doi=true,url=true,isbn=false]{biblatex}
\DeclareLanguageMapping{english}{english}

% Keep index generation under latexmk so out-of-tree builds remain reproducible.
\usepackage{makeidx}
\usepackage{fancyhdr}
\usepackage{emptypage}
\pagestyle{fancy}
\fancyhf{}
\fancyhead[LE]{\small\thepage\quad\leftmark}
\fancyhead[RO]{\small\rightmark\quad\thepage}
\renewcommand{\headrulewidth}{0.3pt}
\fancypagestyle{plain}{\fancyhf{}\fancyfoot[C]{\thepage}\renewcommand{\headrulewidth}{0pt}}

\usepackage{titlesec}
\titleformat{\chapter}[display]
  {\sffamily\bfseries\color{deepblue}}
  {\filleft\Large 第\thechapter 章}
  {0.8ex}
  {\titlerule\vspace{1ex}\Huge}
\titleformat{name=\chapter,numberless}[display]
  {\sffamily\bfseries\color{deepblue}}{}{0pt}{\titlerule\vspace{1ex}\Huge}
\titleformat{\section}{\Large\sffamily\bfseries\color{deepblue}}{\thesection}{0.7em}{}
\titleformat{\subsection}{\large\sffamily\bfseries\color{accent}}{\thesubsection}{0.7em}{}

\usepackage{hyperref}
\hypersetup{
  unicode=true,
  colorlinks=true,
  linkcolor=deepblue,
  citecolor=accent,
  urlcolor=deepblue,
  pdftitle={因果推理深度读本}
}
\usepackage[nameinlink,noabbrev]{cleveref}

\newcommand{\term}[1]{\textbf{#1}\index{#1}}
\newcommand{\eng}[1]{\textit{#1}}
\newcommand{\doop}{\operatorname{do}}
\newcommand{\indep}{\mathrel{\perp\!\!\!\perp}}
\newcommand{\chapterreadings}[1]{\section*{延伸阅读}\addcontentsline{toc}{section}{延伸阅读}\begin{sloppypar}#1\end{sloppypar}}
